\documentclass{beamer}
\usepackage{luatexja}
\usepackage{luatexja}
\renewcommand{\kanjifamilydefault}{\gtdefault}
\usetheme{metropolis}
\title{testSlide}
\author{testAuthor}
\date{\today}

\newcommand{\fanc}{x + 1}
\newcommand{\fancA}[1]{#1 + 1}

\begin{document}
\begin{frame}
    \titlepage
\end{frame}


\begin{frame}
    \section{読みやすいフォント}
    \begin{itemize}
        \item 長文は細いフォント（明朝 or ゴシック）
        \item fontは遊ゴシック、メイリオがおすすめ　場面によって明朝体、英語はセリフ体, time New Romanを使うとよい（英語長文はゴシックだと見ずらい）
        \item Segoe UI、メイリオはおすすめ、ただし、windows以外の環境で見たときに崩れる可能性がある
    \end{itemize}

    メイリオ、time New Roman これだけでも覚えておく

\end{frame}

\begin{frame}
    \section{可能な組み合わせ、特に英数字}
    \begin{itemize}
        \item ゴシック　＋　サンセリフ体
        \item 明朝体＋セリフ体
    \end{itemize}
    これに加えて、英語と文字のサイズを合わせる
\end{frame}


\end{document}